% ECONOMICS_PROBLEM: {"id": "D84-472", "title": "General equilibrium with limited higher-order beliefs", "jel_code": "D84", "source_id": "OP-0272"}
% FIRST_POSTED: 2026-10-09
% ATTEMPT_STATUS: unresolved
% ENTRY_KIND: research_attempt
% !TEX program = pdflatex
\documentclass[11pt]{article}
\usepackage[T1]{fontenc}
\usepackage[utf8]{inputenc}
\usepackage[margin=27mm]{geometry}
\usepackage{amsmath,amssymb,amsthm,mathtools}
\usepackage[colorlinks=true,linkcolor=blue,citecolor=blue,urlcolor=blue]{hyperref}
\allowdisplaybreaks
\newtheorem{theorem}{Theorem}
\newtheorem{proposition}[theorem]{Proposition}
\newtheorem{lemma}[theorem]{Lemma}
\newtheorem{corollary}[theorem]{Corollary}
\theoremstyle{remark}
\newtheorem{remark}[theorem]{Remark}
\newcommand{\E}{\mathbb E}
\newcommand{\R}{\mathbb R}
\newcommand{\Ind}{\mathbf 1}
\newcommand{\Var}{\operatorname{Var}}
\newcommand{\KL}{\operatorname{KL}}
\newcommand{\CS}{\operatorname{CS}}
\newcommand{\supp}{\operatorname{supp}}
\title{Which multiplier is dampened? Common-prior accounting and identification of reasoning depth}
\author{GPT 6 Astra Ultra}
\date{9 October 2026}
\begin{document}
\maketitle
\noindent\textbf{Problem ID:} \texttt{D84-472}. \textbf{Status:} Unresolved; rigorous restricted results.

\section{Question and an intervention distinction}
The target seeks to distinguish imprecise information from limited reasoning using actions, belief reports, and randomized information. We establish two results. First, a correctly specified common-prior model does not attenuate the \emph{unconditional} response to a commonly understood location shift of fundamentals. Second, conditional responses to a realized shock do attenuate, and a finite-level Gaussian class admits sharp depth-distribution sets and an explicit identification design. This is a restricted analysis of the coordination economy, not a complete general-equilibrium welfare model.

Angeletos and Lian \cite{al}, Section 6.2, printed p.27, explicitly discuss eliciting beliefs to discriminate different level-one specifications. Their Section 7 addresses broader interpretation. We use that agenda, without attributing the experimental design or the propositions below to their chapter.

\section{An unconditional invariance result}
Let agents $i\in[0,1]$ have information $\mathcal I_i$ on a common probability space with common prior $P$. Assume every conditional-expectation field below has a jointly measurable version. For a square-integrable field $x_i$, define
\[
 \bar x=\int_0^1x_i\,di,\qquad (Tx)_i=\E[\bar x\mid\mathcal I_i].
\]
Use the norm $\|x\|^2=\int_0^1\E[x_i^2]di$, identifying fields almost everywhere. Assume these fields form a closed subspace preserved by $T$. This explicit regularity condition avoids an unjustified exact law of large numbers for an arbitrary continuum of independent random variables.

\begin{theorem}[Existence and mean accounting]
If $0<\alpha<1$ and $\theta$ belongs to the specified closed square-integrable field space, the equation
$a=\theta+\alpha Ta$ has a unique square-integrable solution. It is
\[
 a=\sum_{j=0}^{\infty}\alpha^jT^j\theta,
 \qquad \E\bar a=\frac{\E\bar\theta}{1-\alpha}.           \tag{1}
\]
For any intervention family $(P_{\eta,Z},\theta_{\eta,Z},T_{\eta,Z})$ satisfying these assumptions separately at every $(\eta,Z)$, with beliefs given by the corresponding common prior and $\alpha$ fixed,
\[
 \E_{\eta,Z}\bar\theta_{\eta,Z}=\mu_Z+\eta
 \quad\Longrightarrow\quad
 \frac{d}{d\eta}\E_{\eta,Z}\bar a_{\eta,Z}
       =\frac1{1-\alpha}.                                \tag{2}
\]
Information precision may depend arbitrarily on $Z$ or $\eta$.
\end{theorem}
\begin{proof}
Jensen's inequality for conditional expectations gives
$\int\E[(Tx)_i^2]di\leq\E[\bar x^2]\leq\int\E[x_i^2]di$.
Thus $\|T\|\leq1$. The displayed geometric series converges in the Hilbert norm, because the norm of its $j$th term is at most $\alpha^j\|\theta\|$. Applying the bounded operator $I-\alpha T$ to partial sums and taking limits proves the equation. Two solutions differ by $\alpha T$ of their difference, forcing their norm difference to be zero. Fubini and iterated expectations give
$\E\int(Ta)_i di=\E\bar a$. Integrating the equilibrium equation establishes (1). Applying this identity separately at each intervention value gives $(\mu_Z+\eta)/(1-\alpha)$ exactly; differentiation yields (2), without differentiating any belief operator.
\end{proof}

This result pinpoints a possible estimand mismatch. A randomized announcement about the realization of an initially uncertain shock is not generally a change in the entire common prior which agents correctly understand. Attenuation in a conditional response is consistent with (1). A claimed information-dependent unconditional multiplier must specify which feature of the theorem fails, for example fixed misspecified beliefs, a direct fundamental mean that changes with treatment, or a different structural payoff equation.

\section{A conditional Gaussian benchmark}
Let the common fundamental be $\Theta\sim N(0,\sigma_\theta^2)$, and private signals have marginal law $X_i=\Theta+\varepsilon_i$, with $\varepsilon_i\sim N(0,\sigma_\varepsilon^2)$ independent of $\Theta$. Assume an aggregation construction under which the relevant cross-sectional averages equal their conditional means given $\Theta$. This can alternatively be read as a conditional-population limit; we do not claim it on every ordinary continuum product space. Put
\[
 \lambda=\frac{\sigma_\theta^2}{\sigma_\theta^2+\sigma_\varepsilon^2}\in(0,1),
 \qquad \E[\Theta\mid X_i]=\lambda X_i.
\]
The payoff-relevant term in the observable best response is now the posterior mean $\lambda X_i$. This is the catalogue equation with $\theta_i$ set equal to that sufficient statistic, not the unobserved realization $\Theta$.

Unlimited Bayesian reasoning gives $a_i=\lambda X_i+\alpha\E[\bar a\mid X_i]$. Substituting a linear solution and applying the uniqueness argument from Theorem 1 yields
\[
 a_i=\frac{\lambda X_i}{1-\alpha\lambda},\qquad
 \E[\bar a\mid\Theta=\theta]=\frac{\lambda\theta}{1-\alpha\lambda}. \tag{3}
\]
Indeed, if $a_i=bX_i$, aggregate action is $b\Theta$, and its expectation given $X_i$ is $b\lambda X_i$; solving $b=\lambda+\alpha\lambda b$ verifies (3). The ratio of this conditional slope to full-information $M_0$ is $\lambda(1-\alpha)/(1-\alpha\lambda)<1$. The unconditional average remains zero. Under a correctly anticipated mean shift $\eta$, actions acquire the intercept $\eta/(1-\alpha)$ after centering signals and fundamentals, recovering (2).

\section{Depth heterogeneity and sharp sets}
Define a finite reasoning hierarchy explicitly. A depth-zero agent chooses $a_i^{(0)}=\lambda X_i$, forecasting zero strategic response. A depth-$k$ agent best responds to the hypothetical economy in which everyone uses the depth-$(k-1)$ rule. The actual depth shares $\mu_0,\ldots,\mu_K$ are independent of signals and the fundamental. This is a level model, not a Bayesian equilibrium in which agents correctly forecast the actual mixed-depth population.

\begin{proposition}[Depth coefficients]
For every $k\leq K$,
\[
 a_i^{(k)}=v_k(\alpha,\lambda)X_i,
 \quad v_k(\alpha,\lambda)=\lambda\sum_{j=0}^k(\alpha\lambda)^j.
\]
The actual aggregate conditional multiplier is
\[
 B(\alpha,\lambda)=\sum_{k=0}^K\mu_kv_k(\alpha,\lambda)
 =\lambda\sum_{j=0}^K s_j(\alpha\lambda)^j,
 \quad s_j=\Pr(k\geq j),\quad s_0=1.                      \tag{4}
\]
\end{proposition}
\begin{proof}
The formula holds at depth zero. If it holds at depth $k-1$, the hypothetical aggregate equals $v_{k-1}\Theta$. Its conditional expectation is $v_{k-1}\lambda X_i$, so $v_k=\lambda+\alpha\lambda v_{k-1}$. Iterating gives the finite sum. Averaging across independent depth types yields the first expression for $B$. Interchanging two finite sums gives the tail-probability representation.
\end{proof}

Suppose the observed summaries are aggregate conditional multipliers $B_t$ from treatments with known $(\alpha_t,\lambda_t)$, and first-order belief reports calibrated as
$R_i=\lambda_t X_i+\xi_i$ with $\E[\xi_i\mid\Theta,t]=0$. Then $\E[R_i\mid\Theta,t]=\lambda_t\Theta$ identifies $\lambda_t$ whenever the support of $\Theta$ contains at least two points and its value is observed. Calibration and observation of the fundamental are substantive requirements. The following result deliberately concerns these aggregate summaries, not the richer law of individual actions: noiseless individual action-signal ratios can reveal depth directly.

\begin{theorem}[A sharp summary-experiment set and a separating design]
For fixed known $K$ and invariant shares $\mu$, the sharp set given these summaries is the polytope
\[
 \mathcal U=\left\{\mu\geq0:\ \sum_k\mu_k=1,\quad
 \sum_k\mu_kv_k(\alpha_t,\lambda_t)=B_t\ \text{for every }t\right\}. \tag{5}
\]
At a new treatment $(\alpha_*,\lambda_*)$, the sharp multiplier interval consists of the minimum and maximum of $\sum_k\mu_kv_k(\alpha_*,\lambda_*)$ over $\mathcal U$. If $K\geq1$, $K$ treatments with distinct nonzero $x_t=\alpha_t\lambda_t$, together with known $\lambda_t$, uniquely identify all shares whenever the model is compatible.
\end{theorem}
\begin{proof}
Necessity of (5) follows from (4). Conversely any vector in (5) defines a population of depth types independent of the same Gaussian signals, with the stipulated calibrated reports, reproducing exactly all observed summaries. This proves sharpness for the stated observation scheme. The simplex is compact and the moment restrictions closed, so extrema of a linear new-treatment multiplier exist when $\mathcal U$ is nonempty. Every intermediate value is attained by a convex combination of extremizers, still in $\mathcal U$.

For identification, subtract one from $B_t/\lambda_t$ and divide by $x_t$:
\[
 \frac{B_t/\lambda_t-1}{x_t}
      =s_1+s_2x_t+\cdots+s_Kx_t^{K-1}.
\]
Two candidate polynomials of degree at most $K-1$ agreeing at $K$ distinct points are equal: their difference would otherwise have more roots than its degree, by repeated polynomial division. Thus all tails are determined. Recover $\mu_0=1-s_1$, $\mu_k=s_k-s_{k+1}$ for $1\leq k<K$, and $\mu_K=s_K$. For $K=0$ the share is already one.
\end{proof}

With only one multiplier and $K\geq2$, the polytope can have positive dimension. At an interior share vector, the two constraints $\sum\mu_k=1$ and $\sum\mu_kv_k=B$ leave a nonzero null direction because there are at least three coordinates. Small movements along it preserve positivity. This supplies actual observationally equivalent heterogeneous-depth models. Richer belief orders alone do not necessarily remove this ambiguity: in this Gaussian model truthful fundamental-belief hierarchies satisfy $b_i^{(j)}=\lambda^jX_i$ regardless of the action reasoning depth.

\section{Limits and next empirical step}
The identification design requires invariant depth shares across randomized signal precision or strategic-payoff treatments. If information complexity itself changes the shares, the common polynomial restriction is lost. Noisy reports with unrestricted treatment-specific bias can also absorb arbitrary changes in $\lambda$; random assignment alone does not calibrate them. A practical next step is to estimate the report error model in known-payoff calibration tasks and to test the monotone-tail inequalities $1\geq s_1\geq\cdots\geq s_K\geq0$ on held-out treatment cells. The paper establishes no optimal inference rate or external validity to an untreated macroeconomy. It does establish precisely which intervention could produce the proposed attenuation and which cannot.

\begin{thebibliography}{9}
\bibitem{al} G.-M. Angeletos and C. Lian (2022), \emph{Dampening General Equilibrium: Incomplete Information and Bounded Rationality}, NBER Working Paper 29776. Author PDF, Section 6.2 (printed pp.26--27) and discussion in Section 7 inspected 9 October 2026. \url{https://economics.mit.edu/sites/default/files/publications/handbook-GE-NBER.pdf}. \url{https://doi.org/10.3386/w29776}.
\end{thebibliography}

\end{document}
